\documentclass{report}
\usepackage{cmt}

\title{\huge{Algorithmische Mathematik I \\ Übungszettel 4}}
\author{\huge{Amy, Rojan}}
\date{2026-11-05}

\begin{document}

\maketitle
\tableofcontents

\chapter{Darstellungen ganzer Zahlen}

\nt{
  Den Übungszettel gibt es
  \href{run:./pset_2026-11-05.pdf}{hier}.
}

\section{Zahlensysteme}

\qs{}{
  \begin{enumerate}
    \item Geben Sie die Darstellung der Zahl $(3324)_5$
      bezüglich der Basis $7$ an.
    \item Wandeln Sie die Oktalzahl $(6272)_8$ in das
      Binärsystem um.
  \end{enumerate}
}

\begin{myproof}
  \textbf{1. Umrechnung von Basis $5$ nach Basis $7$:}

  Zunächst rechnen wir in das Dezimalsystem um:
  \begin{align*}
    (3324)_5
    &=3\cdot5^3+3\cdot5^2+2\cdot5+4\\
    &=375+75+10+4\\
    &=464.
  \end{align*}

  Durch Division mit Rest erhalten wir:
  \begin{align*}
    464&=7\cdot66+2,\\
    66&=7\cdot9+3,\\
    9&=7\cdot1+2,\\
    1&=7\cdot0+1.
  \end{align*}

  Die Reste von unten nach oben ergeben
  \[
    \boxed{(3324)_5=(1232)_7.}
  \]

  \medskip
  \textbf{2. Umrechnung von Basis $8$ nach Basis $2$:}

  Da $8=2^3$, entspricht jede Oktalziffer
  genau drei Binärziffern:
  \begin{align*}
    6&=(110)_2,\\
    2&=(010)_2,\\
    7&=(111)_2,\\
    2&=(010)_2.
  \end{align*}

  Durch Aneinanderhängen folgt
  \[
    \boxed{(6272)_8=(110010111010)_2.}
  \]
\end{myproof}


\section{Zweierkomplement}

\qs{}{
  Sei $z$ die Zweierkomplementdarstellung einer
  negativen Zahl mit $l$ Bits.

  Wie sieht die Darstellung derselben Zahl aus,
  wenn $2l$ Stellen für die Komplementdarstellung
  zur Verfügung stehen?
}

\begin{myproof}
  Sei
  \[
    z=z_{l-1}z_{l-2}\cdots z_0
  \]
  die Zweierkomplementdarstellung einer negativen
  Zahl $x$.

  Da $x<0$, gilt $z_{l-1}=1$.

  Der Zahlenwert ist
  \[
    x=-2^{l-1}
      +\sum_{i=0}^{l-2}z_i2^i.
  \]

  Mit $2l$ Bits erhalten wir dieselbe Zahl,
  indem wir vor die ursprüngliche Darstellung
  $l$ Einsen setzen:
  \[
    \boxed{
      \underbrace{11\cdots1}_{l\text{ Bits}}
      z_{l-1}\cdots z_0.
    }
  \]

  Zum Nachweis betrachten wir den Zahlenwert
  der erweiterten Darstellung:
  \begin{align*}
    x'
    &=-2^{2l-1}
      +\sum_{i=l}^{2l-2}2^i
      +\sum_{i=0}^{l-1}z_i2^i\\
    &=-2^l+(2^l-1)2^l
      \cdot 0
      +\sum_{i=0}^{l-1}z_i2^i.
  \end{align*}

  Einfacher folgt direkt aus
  \[
    \sum_{i=l}^{2l-2}2^i=2^{2l-1}-2^l,
  \]
  dass
  \begin{align*}
    x'
    &=-2^{2l-1}
      +(2^{2l-1}-2^l)
      +\sum_{i=0}^{l-1}z_i2^i\\
    &=-2^l+\sum_{i=0}^{l-1}z_i2^i\\
    &=x.
  \end{align*}

  Die Umwandlung heißt \textbf{Vorzeichenerweiterung}
  (\textit{Sign Extension}).
\end{myproof}


\section{Darstellung zur Basis $-10$}

\dfn{Darstellung zur Basis $-10$}{
  Für Ziffern $a_i\in\{0,\dots,9\}$ definieren wir
  \[
    (a_{l-1}\cdots a_0)_{-10}
    :=\sum_{i=0}^{l-1}a_i(-10)^i.
  \]

  Das Wort $a_{l-1}\cdots a_0$ heißt
  \textbf{Darstellung zur Basis $-10$}, falls
  $a_{l-1}\ne0$ gilt oder die Darstellung
  aus der einzelnen Ziffer $0$ besteht.
}

\qs{}{
  \begin{enumerate}
    \item Schreiben Sie $(19375573910)_{-10}$
      als Dezimalzahl.
    \item Geben Sie eine Darstellung von
      $(9230753)_{10}$ zur Basis $-10$ an.
    \item Zeigen Sie, dass es für jede ganze Zahl
      $x$ eine Darstellung zur Basis $-10$ gibt.
    \item Ist diese Darstellung immer eindeutig?
  \end{enumerate}
}

\begin{myproof}
  \textbf{1. Umrechnung in das Dezimalsystem:}

  Nach Definition gilt
  \begin{align*}
    (19375573910)_{-10}
    &=1(-10)^{10}+9(-10)^9\\
    &\quad+3(-10)^8+7(-10)^7\\
    &\quad+5(-10)^6+5(-10)^5\\
    &\quad+7(-10)^4+3(-10)^3\\
    &\quad+9(-10)^2+1(-10)+0\\
    &=1234567890.
  \end{align*}

  Somit
  \[
    \boxed{(19375573910)_{-10}=1234567890.}
  \]

  \medskip
  \textbf{2. Umrechnung zur Basis $-10$:}

  Wir verwenden wiederholt die Darstellung
  \[
    x=(-10)q+r,
    \qquad r\in\{0,\dots,9\}.
  \]

  Dadurch erhalten wir:
  \begin{align*}
    9230753&=(-10)(-923075)+3,\\
    -923075&=(-10)(92308)+5,\\
    92308&=(-10)(-9230)+8,\\
    -9230&=(-10)(923)+0,\\
    923&=(-10)(-92)+3,\\
    -92&=(-10)(10)+8,\\
    10&=(-10)(-1)+0,\\
    -1&=(-10)(1)+9,\\
    1&=(-10)(0)+1.
  \end{align*}

  Die Reste ergeben von unten nach oben
  \[
    \boxed{(9230753)_{10}=(190830853)_{-10}.}
  \]

  \medskip
  \textbf{3. Existenz:}

  Wir zeigen die Existenz einer Darstellung
  durch Induktion über $|x|$.

  \textbf{Induktionsanfang:}

  Für $x=0$ gilt die Darstellung $(0)_{-10}$.
  Für $x=1$ gilt $(1)_{-10}$, und für $x=-1$ gilt
  \[
    -1=1\cdot(-10)+9=(19)_{-10}.
  \]

  \textbf{Induktionsschritt:}

  Sei nun $|x|\ge2$ und angenommen, alle
  ganzen Zahlen mit kleinerem Betrag besitzen
  eine Darstellung zur Basis $-10$.

  Nach Division mit Rest existieren eindeutig
  bestimmte ganze Zahlen $q$ und
  $r\in\{0,\dots,9\}$ mit
  \[
    x=(-10)q+r.
  \]

  Dabei gilt
  \[
    q=\frac{r-x}{10}.
  \]

  Für $|x|\ge2$ ist $|q|<|x|$.
  Nach Induktionsvoraussetzung besitzt $q$
  somit eine Darstellung zur Basis $-10$.

  Sei
  \[
    q=\sum_{i=0}^{l-1}a_i(-10)^i.
  \]

  Dann folgt
  \begin{align*}
    x
    &=(-10)q+r\\
    &=(-10)\sum_{i=0}^{l-1}a_i(-10)^i+r\\
    &=r+\sum_{i=1}^{l}a_{i-1}(-10)^i.
  \end{align*}

  Also besitzt auch $x$ eine Darstellung
  zur Basis $-10$.

  Damit existiert eine solche Darstellung
  für jede ganze Zahl $x$.

  \medskip
  \textbf{4. Eindeutigkeit:}

  Angenommen, $x$ besitze zwei Darstellungen
  zur Basis $-10$.

  Die letzte Ziffer ist jeweils eindeutig durch
  \[
    a_0\equiv x\pmod{10},
    \qquad a_0\in\{0,\dots,9\}
  \]
  bestimmt.

  Daher müssen beide Darstellungen dieselbe
  letzte Ziffer besitzen.

  Nach Abziehen von $a_0$ und Division
  durch $-10$ erhalten wir
  \[
    \frac{x-a_0}{-10}.
  \]

  Dieses Argument lässt sich wiederholt
  anwenden, sodass auch alle weiteren
  Ziffern übereinstimmen müssen.

  Da führende Nullen ausgeschlossen sind,
  ist auch die Länge eindeutig.

  Somit ist die Darstellung zur Basis $-10$
  für jede ganze Zahl eindeutig.
\end{myproof}


\section{Quersummen}

\dfn{Quersumme}{
  Sei
  \[
    n=\sum_{i=0}^{l-1}z_i10^i
  \]
  die Dezimaldarstellung einer natürlichen Zahl.

  Dann heißt
  \[
    q(n):=\sum_{i=0}^{l-1}z_i
  \]
  die \textbf{Quersumme} von $n$.

  Für $t\in\bbN$ definieren wir rekursiv
  \[
    q_1(n):=q(n),
    \qquad
    q_{t+1}(n):=q(q_t(n)).
  \]

  Ferner definieren wir
  \[
    q^*(n):=q_{t_0}(n),
  \]
  wobei $t_0$ der kleinste natürliche Index
  mit $q_{t_0}(n)\le9$ ist.
}

\qs{}{
  Seien $a,b\in\bbN$. Zeigen Sie:
  \[
    \boxed{
      q^*(q^*(a)\cdot q^*(b))
      =q^*(a\cdot b).
    }
  \]
}

\begin{myproof}
  Wir verwenden die Kongruenz
  \[
    10\equiv1\pmod9.
  \]

  Sei
  \[
    n=\sum_{i=0}^{l-1}z_i10^i.
  \]
  Dann folgt
  \begin{align*}
    n
    &\equiv\sum_{i=0}^{l-1}z_i
      \pmod9\\
    &\equiv q(n)\pmod9.
  \end{align*}

  Durch wiederholte Anwendung erhalten wir
  \[
    q^*(n)\equiv n\pmod9.
  \]

  Außerdem gilt für alle $n\in\bbN$:
  \[
    q^*(n)\in\{1,\dots,9\}.
  \]

  Damit ist $q^*(n)$ die eindeutig bestimmte
  Zahl aus $\{1,\dots,9\}$, die zu $n$
  modulo $9$ kongruent ist.

  Für $a,b\in\bbN$ folgt
  \begin{align*}
    q^*(a)\cdot q^*(b)
    &\equiv a\cdot b\pmod9.
  \end{align*}

  Also gilt ebenso
  \[
    q^*(q^*(a)\cdot q^*(b))
    \equiv a\cdot b\pmod9.
  \]

  Da beide Werte
  \[
    q^*(q^*(a)\cdot q^*(b))
    \quad\text{und}\quad
    q^*(a\cdot b)
  \]
  in $\{1,\dots,9\}$ liegen und modulo $9$
  kongruent sind, müssen sie gleich sein.

  Somit gilt
  \[
    \boxed{
      q^*(q^*(a)\cdot q^*(b))
      =q^*(a\cdot b).
    }
  \]
\end{myproof}


\section{Rechnen mit rationalen Zahlen}

\qs{}{
  Zeigen Sie, dass die Addition und Multiplikation
  zweier rationaler Zahlen, deren Zähler und
  Nenner jeweils höchstens $l$ Stellen haben,
  in $O(l^2)$ Rechenschritten durchgeführt werden
  können.
}

\begin{myproof}
  Seien
  \[
    x=\frac ab,
    \qquad
    y=\frac cd,
    \qquad b,d\ne0,
  \]
  wobei $a,b,c,d$ jeweils höchstens $l$
  Dezimalstellen besitzen.

  Wir verwenden die schriftlichen
  Rechenverfahren.

  \medskip
  \textbf{Addition:}

  Es gilt
  \[
    \frac ab+\frac cd
    =\frac{ad+bc}{bd}.
  \]

  Dafür müssen wir drei Multiplikationen
  und eine Addition durchführen.

  Die schriftliche Multiplikation zweier
  Zahlen mit höchstens $l$ Stellen benötigt
  $O(l^2)$ Ziffernoperationen.

  Somit benötigen
  \[
    ad,\qquad bc,\qquad bd
  \]
  jeweils $O(l^2)$ Schritte.

  Die Addition von $ad$ und $bc$ benötigt
  höchstens $O(l)$ Schritte, da beide Produkte
  höchstens $2l$ Stellen besitzen.

  Daher gilt insgesamt
  \[
    T_{\mathrm{Add}}(l)
    =3\cdot O(l^2)+O(l)
    =O(l^2).
  \]

  \medskip
  \textbf{Multiplikation:}

  Es gilt
  \[
    \frac ab\cdot\frac cd
    =\frac{ac}{bd}.
  \]

  Hier sind lediglich zwei schriftliche
  Multiplikationen nötig.

  Beide benötigen jeweils $O(l^2)$ Schritte.
  Somit folgt
  \[
    T_{\mathrm{Mult}}(l)
    =2\cdot O(l^2)
    =O(l^2).
  \]

  Also können Addition und Multiplikation
  rationaler Zahlen mit höchstens $l$ Stellen
  in Zähler und Nenner jeweils in
  \[
    \boxed{O(l^2)}
  \]
  Ziffernoperationen durchgeführt werden.
\end{myproof}

\nt{
  Dabei wird das Ergebnis als Bruch dargestellt.
  Ein anschließendes vollständiges Kürzen des
  Bruchs ist für diese Abschätzung nicht erforderlich.
}


\section{Programmierteil}

\qs{}{
  Schreiben Sie
  \begin{enumerate}
    \item das Programm \texttt{prime.cpp}
      (einfacher Primzahltest) so um, dass
      anstelle der \texttt{for}-Schleife eine
      \texttt{while}-Schleife benutzt wird;

    \item das Programm \texttt{collatz.cpp}
      (Collatz-Folge) so um, dass anstelle der
      \texttt{while}-Schleife eine
      \texttt{for}-Schleife benutzt wird.
  \end{enumerate}

  Es sind nur die in der Vorlesung behandelten
  C++-Sprachelemente erlaubt.
}

\subsection{Primzahltest mit while-Schleife}
\codefile{prime.cpp}{C++}

\subsection{Collatz-Folge mit for-Schleife}
\codefile{collatz.cpp}{C++}
\end{document}
